% !TeX spellcheck = ru_RU
% !TEX root = vkr.tex

\section{Заключение} \label{sec:conclusion}

В рамках настоящей работы проведено экспериментальное исследование влияния дообучения большой языковой модели YandexGPT~5.1 методом LoRA~\cite{hu2022lora} на качество классификации инструментов в AI-агенте, функционирующем в домене недвижимости. Исследование выполнено на двух независимых бенчмарках общим объёмом 750~примеров (500 смешанных и 250 реальных) и охватывает 12~моделей с вариацией типа и объёма обучающих данных.

Основные задачи, поставленные в начале исследования, были решены в полном объёме:
\begin{enumerate}
    \item Формализована задача классификации инструментов как задача бинарной классификации с классами GEO (\texttt{ResolveEntities}) и TAG (\texttt{SearchTags}); сформированы два независимых бенчмарка для раздельной и совместной оценки моделей.
    \item Разработан модульный программный пайплайн, обеспечивающий полный цикл~--- от подготовки обучающих данных и запуска дообучения через API YandexGPT до автоматизированного inference и вычисления метрик с визуализацией результатов.
    \item Подготовлено 12~конфигураций моделей (1 базовая и 11 дообученных) с систематической вариацией источника данных (синтетические, реальные, смешанные), классового баланса и объёма выборки (от~50 до~${\sim}1500$~примеров), и проведено 24 цикла оценки на двух бенчмарках.
    \item Выполнен строгий статистический анализ результатов с применением критерия Кохрена~$Q$ для проверки глобальной гетерогенности и попарного критерия Макнемара с поправкой Бонферрони для множественных сравнений.
\end{enumerate}

Главные научные результаты работы заключаются в следующем:
\begin{itemize}
    \item \textbf{Поддержана рабочая гипотеза о стилистической природе выбора инструмента.} Дообучение методом LoRA снижает error rate с $18{,}2\%$ до $0{,}6\%$ (более чем в~$30$~раз; $91 \to 3$ ошибочных ответов из~$500$), а Macro~$F_1$ возрастает с $0{,}526$ до $0{,}996$. Различия статистически значимы (раздел~\ref{subsec:significance}). Результат согласуется с гипотезой поверхностного выравнивания~\cite{zhou2023lima}: небольшой объём примеров корректного формата достаточен для радикального улучшения качества. Это эмпирическая поддержка аналогии, а не строгое доказательство самой SAH применительно к tool selection~--- вопрос о том, действительно ли дообучение затрагивает лишь <<форму>>, а не <<знания>> модели в этом контексте, выходит за рамки настоящей работы.

    \item \textbf{Установлена оптимальная стратегия формирования обучающей выборки.} Модель на смешанном датасете (250~реальных + 250~синтетических примеров) достигает наилучших результатов и демонстрирует минимальную деградацию при переходе на реальные данные. Увеличение объёма синтетических данных сверх~500 примеров приводит к ухудшению, что свидетельствует о переобучении на артефакты генерации. Качество обучающих данных оказывается важнее их объёма (раздел~\ref{subsec:data_type}).

    \item \textbf{Выявлен эффект консервации ошибок.} $73{,}6\%$~примеров бенчмарка классифицируются одинаково всеми моделями, что указывает на устойчивое ядро <<лёгких>> и <<трудных>> запросов. Ошибки воспроизводятся на обоих бенчмарках, что свидетельствует об их детерминированной природе~--- они обусловлены систематическими слабостями модели на определённых типах запросов, а не стохастикой генерации.
\end{itemize}

Практическая значимость работы определяется следующим. Во-первых, разработанный программный пайплайн (описанный в разделе~\ref{sec:implementation} и предшествующей работе~\cite{goroshko2025otchet}) может быть использован для дообучения и систематической оценки моделей YandexGPT в произвольных задачах классификации. Во-вторых, сформулированы конкретные рекомендации по формированию обучающих данных: использование смешанных выборок (реальные + синтетические) предпочтительнее наращивания объёма чисто синтетических примеров, что позволяет снизить затраты на ручную разметку без потери качества. В-третьих, показано, что метод LoRA обеспечивает переход от неприемлемого уровня качества к практически безошибочной классификации (раздел~\ref{subsec:main_results}). Ограничения, накладываемые на интерпретацию полученных результатов, систематизированы в разделе~\ref{sec:threats}.

\subsection{Направления дальнейших исследований}

Проведённое исследование открывает ряд перспективных направлений для продолжения работы:
\begin{itemize}
    \item \textbf{Переход к multi-label классификации.} В текущей постановке каждый запрос относится ровно к одному классу. Однако на практике встречаются запросы, которые выглядят географически, но на самом деле описывают характеристики инфраструктуры~--- например, <<рядом с парком>>, <<у озера>>, <<недалеко от леса>>. Такие запросы классифицируются как TAG, хотя содержат пространственные указания. Расширение модели до multi-label классификации позволит корректно обрабатывать подобные пограничные случаи.

    \item \textbf{Исследование альтернативных PEFT-методов.} В данной работе использован исключительно метод LoRA. Сравнение с другими подходами к параметрически эффективному дообучению~--- QLoRA, Prompt Tuning~\cite{lester2021power}, Adapter Tuning~--- позволит определить, является ли LoRA оптимальным выбором для задачи tool selection.

    \item \textbf{Обобщение на другие LLM.} Все эксперименты проведены на семействе моделей YandexGPT~5.1. Проверка воспроизводимости результатов на моделях иных архитектур и провайдеров позволит установить, в какой мере полученные закономерности являются универсальными, а в какой~--- специфичными для конкретной модели.
\end{itemize}

Таким образом, настоящая работа вносит вклад в понимание механизмов адаптации больших языковых моделей к задаче выбора инструментов и демонстрирует, что дообучение методом LoRA на относительно небольшой, но качественно сформированной обучающей выборке является эффективной стратегией повышения качества классификации в AI-агентах.
